% BHU PYQ -- Manually transcribed & error-corrected
% Paper: GLB-501: Physical and Structural Geology
% Exam: B.Sc. (Hons.) Semester V Examination, 2022-23

\documentclass[11pt,a4paper]{article}
\usepackage[a4paper,top=2.5cm,bottom=2.5cm,left=2.8cm,right=2.8cm,headheight=14pt]{geometry}
\usepackage{amsmath,amssymb}
\usepackage[T1]{fontenc}
\usepackage[utf8]{inputenc}
\usepackage{lmodern,microtype,fancyhdr,enumitem,hyperref,lastpage,parskip}

\hypersetup{colorlinks=true,urlcolor=black,linkcolor=black,
  pdftitle={GLB-501 Physical and Structural Geology -- BHU B.Sc. Sem V 2022-23}}
\pagestyle{fancy}\fancyhf{}
\renewcommand{\headrulewidth}{0.4pt}\renewcommand{\footrulewidth}{0.4pt}
\fancyhead[L]{\small   \textit{Banaras Hindu University}}
\fancyhead[R]{\small   \textit{GLB-501: Physical \& Structural Geology}}
\fancyfoot[C]{\small Page  \thepage\ of \pageref{LastPage}}


\newlist{parts}{enumerate}{2}
\setlist[parts,1]{label=(\alph*),leftmargin=*,itemsep=5pt,topsep=3pt}

\newcommand{\pts}[1]{\hfill   \textit{[#1]}}


\begin{document}

\begin{center}
  {\large\bfseries BANARAS HINDU UNIVERSITY}\\[2pt]
  {
\normalsize B.Sc. (Hons.) --- Semester V Examination, 2022-23}\\[6pt]
  \rule{0.8\linewidth}{0.5pt}\\[6pt]
  {\large\bfseries Paper: GLB-501 --- Physical and Structural Geology}\\[4pt]
  {
\normalsize Time: 3 Hours \hfill Maximum Marks: 70}\\[2pt]
  {\small Ref. No.: 052735}
\end{center}
\rule{\linewidth}{0.4pt}
\smallskip



\noindent   \textit{Instructions: Answer   \textbf{five} questions in all, selecting at least   \textbf{two} questions from each of Sections A and B. Section-C is compulsory. All questions carry equal marks.}
\bigskip

\bigskip


\noindent {\large\bfseries SECTION --- A}
\rule{\linewidth}{0.4pt}\smallskip




\noindent   \textbf{Question 1.} Describe the erosional and depositional features formed by glaciers. \pts{14}

\medskip


\noindent   \textbf{Question 2.} Briefly discuss the following: \pts{$14 = 4+5+5$}

\begin{parts}
  \item Landscape Plain
  \item Wegener's Continental Drift Theory
  \item Drainage patterns
\end{parts}

\medskip


\noindent   \textbf{Question 3.} Write a short note on the following: \pts{$14 = 4+5+5$}

\begin{parts}
  \item Mid-oceanic ridges with suitable examples
  \item Isostasy, proposed by Pratt and Airy
  \item Define landslide and discuss the mitigation of landslides
\end{parts}

\bigskip


\noindent {\large\bfseries SECTION --- B}
\rule{\linewidth}{0.4pt}\smallskip




\noindent   \textbf{Question 4.} Define Joint and different joint terminologies with a neat diagram. Discuss the classification of joints based on spatial relationship and geometry with neat diagrams. \pts{14}

\medskip


\noindent   \textbf{Question 5.} What are the fundamental possibilities of stress regimes and stress orientations that produce three types of faults? Discuss with neat diagrams. \pts{14}

\medskip


\noindent   \textbf{Question 6.} Describe the difference between the following: \pts{$14 = 4+4+4+2$}

\begin{parts}
  \item Window and Klippe \pts{4}
  \item Listric and Detachment fault \pts{4}
  \item Reclined and Recumbent folds \pts{4}
  \item S-fabrics and L-fabrics \pts{2}
\end{parts}

\medskip


\noindent   \textbf{Question 7.} Discuss different parameters required to understand Ramsay's geometrical classification of folds with a neat diagram. \pts{14}

\bigskip


\noindent {\large\bfseries SECTION --- C}
\rule{\linewidth}{0.4pt}\smallskip




\noindent   \textbf{Question 8.} Discuss each of the following in two or three sentences. Give suitable diagrams wherever necessary: \pts{$14 = 3+3+3+3+2$}

\begin{parts}
  \item Spherical and stereographic projection of a dipping plane \pts{3}
  \item Nonconformity \pts{3}
  \item Onshore, Backshore, Foreshore, and Offshore \pts{3}
  \item Epirogeny and orogeny \pts{3}
  \item Strain distribution on a profile of a buckling fold \pts{2}
\end{parts}

\vfill\rule{\linewidth}{0.4pt}

\begin{center}\small
\href{https://bhu-exam.vercel.app/}{Click here to visit our website}\\[4pt]
\href{https://me-aryan.vercel.app/}{Click here to visit our contributor}\\[4pt]
\href{https://quashan-me.vercel.app/}{Click here to visit our contributor}
\end{center}
\end{document}
